\subsection{Locally nilpotent upper-fiber ideals}
We use ordinary Levitzki-radical facts; see
\cite[p.~53 and Corollary~1, p.~57]{Rich}. The associative finite-corner
argument needed here is included, without the additional hypotheses
used for alternative rings in that source.

The global semiring Levitzki radical of Olson and Jenkins
\cite[Definitions~1--2, p.~91]{OlsonJenkins} uses products vanishing at
the global zero and Bourne ideal quotients. The present construction
uses quotients of the coefficient groups and retains the support.
For example, the global radical of $R_0$ is $\{\tau\}$, since no
old-ring element has a power equal to $\tau$, whereas our old
coefficient fiber is $\operatorname{Lev}(R)$ and can be nonzero.
Thus the two notions of radical are distinct.

\begin{proof}[Proof of Theorem~\ref{thm:upper-levitzki-quotient}]
First, locally nilpotent rings are closed under extensions by ideals.
If $I\lhd B$ and both $I$ and $B/I$ are locally nilpotent, take a finite
set $X\subset B$. All words of some length $a$ in $X$ lie in $I$.
There are finitely many such words, and they generate a nilpotent
subring, of some exponent $b$. Every word of length $ab$ in $X$
vanishes. This also proves that sums of locally nilpotent ideals are
locally nilpotent: use finite sums and then the finite number of
summands needed for any finite set. Their sum is
$\operatorname{Lev}(B)$, and its quotient has no nonzero locally
nilpotent ideal, by the same extension argument.

The context tests define additive subgroups $H_q$. Transition
stability follows by transporting a test to its later source;
two-sided multiplication stability follows by replacing
$u(axb)v$ by $(ua)x(bv)$. At an upper idempotent $e$, the unit test
gives $H_e\subseteq\operatorname{Lev}(F(e))$. This is an ideal of
$F(e)$ and is locally nilpotent. Conversely, any coefficient ideal
locally nilpotent at all these fibers passes every context test, which
proves maximality. The zero-test ideal $K^{\mathrm{id}}$ is contained
in $H$. For a quotient coefficient ideal with the stated property,
its inverse-image fiber at $e$ is an extension of $H_e$ by a locally
nilpotent ideal. It is locally nilpotent, so maximality makes that
inverse image $H$. This proves the residual assertion. Pair
isomorphisms preserve all tests.

We next prove the ordinary finite-corner fact. Let $B$ be unital with
orthogonal idempotents $c_1,\ldots,c_n$ summing to $1$, and let
$J\lhd B$ have all $c_iJc_i$ locally nilpotent. Then $J$ is locally
nilpotent. Induct on $n$. For a finite subset of $J$, include all its
Peirce components $c_ixc_j$. Put $c=c_1$ and $f=1-c$. Induction
makes $fJf$ locally nilpotent, so words in the finitely many
$f$-to-$f$ generators vanish at some length $a$. In a nonzero block
word, an excursion from $c$ through $f$ and back therefore has the
form
\[
 x_{cf}y_1\cdots y_sx_{fc},\qquad 0\leq s<a.
\]
These finitely many coefficients, together with the direct
$c$-to-$c$ generators, lie in $cJc$ and generate a nilpotent subring,
of some exponent $b$. A nonzero path has fewer than $b$ complete
$c$-to-$c$ pieces. Its initial and terminal $f$ segments and every
excursion have bounded length. Thus sufficiently long words in the
finite component alphabet vanish. This proves the corner fact for
arbitrary associative rings.

Fix $E^2=E$ and put $D=E\vee I_n$, entrywise. Then $D^2=D$,
$D_{ii}=e_i=E_{ii}\vee1\in U(Q)$, and the off-diagonal entries are
unchanged. The ordinary ring $B=A_D(T)$ is unital: its identity has
the unit of $F(e_i)$ on the $i$th diagonal and zero coefficients
elsewhere. Indeed,
\[
 e_iE_{ij}=E_{ij}=E_{ij}e_j.
\]
Its coordinate idempotents $c_i$ have corners
$c_iBc_i\simeq F(e_i)$. These are identities in one fixed ordinary
ring fiber, not a change of the semiring identity support.

Entrywise transition gives a ring map
$\phi:A_E(T)\to A_D(T)$. Its kernel has only diagonal entries, in
$\ker(F(E_{ii})\to F(e_i))$. Such an entry annihilates every entry
in its row or column: the product factors through its zero image in
$F(e_i)$ by the displayed support identities. Therefore
\[
 (\ker\phi)A_E=0=A_E(\ker\phi).
\]
The ideal $N_D\lhd B$ with entries $H_{D_{ij}}$ has corner ideals
$H_{e_i}$, all locally nilpotent. The finite-corner fact makes $N_D$
locally nilpotent. The kernel $N_E$ of $A_E(T)\to A_E(T/H)$ maps
into $N_D$ under $\phi$, and its remaining kernel has square zero.
The extension argument therefore makes $N_E$ locally nilpotent.

Lift a quotient idempotent at its exact support $E$ to $a\in A_E$.
Its defect $d=a^2-a\in N_E$ is nilpotent. The correction
$a\mapsto3a^2-2a^3$ changes $d$ to $d^2(4d-3)$, so iteration gives
an exact idempotent without a constant term or change of support.
For an isomorphism, lift its rectangular inverse intertwiners and cut
by lifted idempotents $P,Q$. The resulting composites are
$XY=P+\delta$ and $YX=Q+\epsilon$. Each error is nilpotent, hence
both composites are invertible by finite geometric series in their
corners. As before, $Y'=Y(XY)^{-1}=(YX)^{-1}Y$ is an exact inverse.
This proves existence and uniqueness of lifted projective classes.
Extension preserves direct sums, proving the $V$ and $K_0$ assertions.

Finally let $S=M_m(T)$. The entrywise family $M_mH(T)$ is a
coefficient ideal. At every upper idempotent support matrix its fiber
is locally nilpotent by what was just proved, so maximality gives
$M_mH(T)\subseteq H(S)$. Conversely take $X=(x_{ij})\in H(S)_P$.
For a scalar test $u\in F(p)$, $v\in F(r)$,
$pP_{ij}r\leq e\in U(Q)$, put $u$ in position $(1,i)$ of $U$ and
$v$ in position $(j,1)$ of $V$, with global semiring zeros elsewhere.
Transport $UXV$ to the upper idempotent support
$D=\operatorname{diag}(e,1,\ldots,1)$. Its coefficient belongs to
$\operatorname{Lev}(A_D(T))$. The first ordinary Peirce corner is
$F(e)$, and the corner of a locally nilpotent ideal is locally
nilpotent. Thus the selected scalar coefficient
$t_{pP_{ij}r,e}(ux_{ij}v)$ belongs to $\operatorname{Lev}(F(e))$.
Every entry passes all tests, giving $H(S)\subseteq M_mH(T)$.
This argument allows $F(0)\ne0$; it does not identify $A_D$ with a
direct product. Entrywise quotient gives the asserted matrix
isomorphism.
\end{proof}

\begin{example}[ordinary boundary and the residual condition]
\label{ex:levitzki-boundary}
For an ordinary ring $R$, its one-point support gives
$H=\operatorname{Lev}(R)$ and $K^{\mathrm{id}}=0$. Dual numbers over
a field give strict containment. The rings $k[t]/(t^a)$ have
$1\times1$ fiber kernels $(t)$ of exponent $a$, so no nilpotence
bound depending only on matrix size exists.

Individual upper fibers need not be radical-free after the reduction.
Take $Q=\{\bot<1<e\}$, with maximum for join and for nonbottom
multiplication. Put $F(\bot)=0$, $F(1)=k[t]/(t^2)$ and
$F(e)=M_2(k)$, with the injective unital transition
$a+bt\mapsto aI+bE_{12}$. Products use this map and ordinary ring
products, so the pair axioms hold. Since $M_2(k)$ has no nonzero
locally nilpotent ideal, $H_e=0$; injectivity of the transition then
gives $H_1=0$. Thus $H=0$, although
$\operatorname{Lev}(F(1))=(t)\ne0$. The residual property is exactly
the absence of a compatible coefficient ideal of the specified kind.
\end{example}

The zero-detection ideal alone need not commute with matrices. In
Example~\ref{ex:upper-idempotent-strict}, the support matrix with
integer-zero diagonal, upper off-diagonal $1\in\mathbb N$ and lower
off-diagonal $\bot$ is an upper idempotent. The entrywise
$K^{\mathrm{id}}$ has a nonzero upper off-diagonal group there,
whereas $K^{\mathrm{id}}(M_2(T))$ vanishes on the entire fiber by its
unit test. The locally nilpotent quotient includes this ordinary
triangular radical and has the matrix compatibility just proved.

